As a field, stellar modelling is now seven decades old. Yet, there remain challenges associated with the numerical method; chief among them the treatment of secondaries, and modelling the core-collapse supernova and its aftermath. In this thesis, we present self-consistent calculations for three major phases of stellar evolution. We display a numerical example of a system in which helium-rich material is transferred between the constituent stars, and compare this to one where the composition of the accreted mass is fixed to the companion's surface composition. We find that this is results in accretors which are hotter and at times more luminous than when material is assumed to reflect the surface composition of the secondary. We present a method for determining the goodness-of-fit of a natal supernova kick, and display this method for a kick linear in the ejecta-to-remnant mass ratio $Q$: $v_\text{kick}=\alpha Q + \beta$, with free parameters $\alpha$ and $\beta$. We find that the parameters $\alpha=\SI[parse-numbers=false]{115^{+40}_{-55}}{\kilo\meter\per\second}$ and $\beta=\SI[parse-numbers=false]{15^{+10}_{-15}}{\kilo\meter\per\second}$ are able to replicate the merger rate of observed galactic double neutron stars, period eccentricity distributions, and isolated pulsar velocities, and produce ultra-stripped supernovae. We present a small grid of overcontact binary stars, where both stars have been evolved in detail. We use the BPASS v2.2 models in order to constrain the initial zero-age main sequence parameters of six observed O+O overcontact systems. We use these parameters as inputs to a variant of the Cambridge \codename{stars} code, and produce matches to one of the binaries, with candidate matches for two more. We discuss avenues for future research, including modifications to the mass-loss prescriptions already in use within the code, as well as population studies of overcontact binaries.